\documentclass[11pt]{article}
\usepackage[margin=1in]{geometry}
\usepackage{amsmath,amssymb,amsthm}
\usepackage{graphicx}
\usepackage[hidelinks]{hyperref}

\newtheorem{theorem}{Theorem}
\newtheorem{lemma}{Lemma}
\newtheorem{corollary}{Corollary}

\newcommand{\R}{\mathbb R}
\newcommand{\HH}{\mathbb H}
\newcommand{\dd}{\,d}
\newcommand{\norm}[1]{\left\lVert #1\right\rVert}
\newcommand{\abs}[1]{\left\lvert #1\right\rvert}

\title{A Step-Two Carnot Subcase of Ryoo's Endpoint \(L^1\) Question}
\author{Run \texttt{fa\_banach\_001}, agent \texttt{agent\_lane\_16}}
\date{June 15, 2026}

\begin{document}
\maketitle

\section*{Status}

This packet records a likely valid partial solution to Question 38 / Question 42 in Ryoo's paper
\emph{Quantitative nonembeddability of groups of polynomial growth into uniformly convex spaces}
\cite{Ryoo2025}. The result proves that every nonabelian step-two Carnot group has a finite endpoint
\(L^1\) exponent, namely \(q=4\). It does not claim optimality in every step-two group.

\section*{Source question}

Ryoo asks for which exponents \(q\geq 1\) every compactly supported smooth
\(f:G\to \R\) satisfies
\begin{equation}\label{eq:source-question}
\left(\int_0^\infty
\left(\int_G \abs{f(h\exp(tv))-f(h)}\dd\mu(h)\right)^q
\frac{\dd t}{t^{1+q/s}}\right)^{1/q}
\lesssim_{G,q}
\int_G \norm{\nabla f(h)}_{\ell_2^k}\dd\mu(h),
\end{equation}
where \(G\) is a nonabelian simply connected nilpotent Lie group of step \(s\) and \(v\) is in the
last nonzero lower-central term. The source notes that this is equivalent, by coarea, to a vertical
perimeter versus horizontal perimeter inequality.

\begin{center}
\includegraphics[width=0.94\linewidth]{figures/open_problem_crop.png}
\end{center}

\section*{Claim proved}

Let \(G\) be a step-two Carnot group with Lie algebra
\(\mathfrak g=V_1\oplus V_2\), where \(V_2=[V_1,V_1]\) is central. Equip \(V_1\) with the fixed
Carnot inner product and write \(\nabla_H f\) for the horizontal gradient. For \(z\in V_2\), define
\[
A_zf(t)=\int_G \abs{f(g\exp(tz))-f(g)}\dd g,\qquad t>0.
\]

\begin{theorem}\label{thm:main}
For every \(v\in V_2\) there is a constant \(C_{G,v}<\infty\) such that every
\(f\in C_c^\infty(G)\) satisfies
\begin{equation}\label{eq:q4}
\left(\int_0^\infty \frac{A_vf(t)^4}{t^3}\dd t\right)^{1/4}
\leq C_{G,v}\int_G \abs{\nabla_H f(g)}\dd g .
\end{equation}
Consequently, after normalizing \(d_G(\exp v,e)=1\), Question \eqref{eq:source-question} has an
affirmative answer with \(q=4\) for every step-two Carnot group.
\end{theorem}

\section*{Proof intuition}

Every central direction in a step-two Carnot group is a finite sum of elementary commutators
\([a,b]\) with \(a,b\in V_1\). Each nonzero elementary commutator generates a three-dimensional
Heisenberg subgroup. Naor and Young proved the exact \(L^4\) endpoint vertical-horizontal
inequality in \(\HH^3\) \cite{NaorYoung2022}. Applying that theorem on each left coset of the
Heisenberg subgroup controls translation in the direction \([a,b]\) by the ambient horizontal
gradient. Since all \(V_2\)-directions are central in a step-two group, translating by a sum of
commutators can be telescoped into translations by the individual commutators.

\section*{Supporting theorem used}

Naor and Young's Theorem 1.1 says, in functional form, that in the three-dimensional Heisenberg
group
\[
\left(\int_0^\infty
\left(\int_{\HH^3}\frac{\abs{F(h\exp(tZ))-F(h)}}{\sqrt t}\dd h\right)^4
\frac{\dd t}{t}\right)^{1/4}
\lesssim
\int_{\HH^3} \bigl(\abs{XF(h)}+\abs{YF(h)}\bigr)\dd h .
\]
Equivalently,
\[
\left(\int_0^\infty
\frac{\left(\int_{\HH^3}\abs{F(h\exp(tZ))-F(h)}\dd h\right)^4}{t^3}\dd t
\right)^{1/4}
\lesssim
\int_{\HH^3} \bigl(\abs{XF(h)}+\abs{YF(h)}\bigr)\dd h .
\]
The relevant theorem statement is reproduced below for context.

\begin{center}
\includegraphics[width=0.90\linewidth]{figures/supporting_h3_theorem_crop.png}
\end{center}

\section*{Proof}

We first transfer the \(\HH^3\) theorem to elementary bracket directions in \(G\).

\begin{lemma}\label{lem:elementary}
Let \(a,b\in V_1\) and \(z=[a,b]\neq 0\). Then there is a constant \(C_{a,b,G}\) such that
\[
\left(\int_0^\infty \frac{A_zf(t)^4}{t^3}\dd t\right)^{1/4}
\leq
C_{a,b,G}\int_G \bigl(\abs{a f(g)}+\abs{b f(g)}\bigr)\dd g
\leq
C_{a,b,G}'\int_G \abs{\nabla_H f(g)}\dd g
\]
for every \(f\in C_c^\infty(G)\). Here \(a f\) and \(b f\) denote the left-invariant derivatives
in the horizontal directions \(a\) and \(b\).
\end{lemma}

\begin{proof}
Let \(H\leq G\) be the connected subgroup with Lie algebra
\(\mathfrak h=\operatorname{span}\{a,b,z\}\). Since \(z=[a,b]\neq 0\) and \(G\) has step two,
\(\mathfrak h\) is a three-dimensional Heisenberg Lie algebra, up to a linear rescaling of the
horizontal basis and central coordinate.

Because \(G\) and \(H\) are unimodular nilpotent Lie groups, Weil integration gives a quotient
measure on \(G/H\) such that
\[
\int_G \Phi(g)\dd g=\int_{G/H}\int_H \Phi(xu)\dd u\,\dd \bar x
\]
for compactly supported measurable \(\Phi\). For each coset representative \(x\), apply the
Naor--Young \(\HH^3\) inequality to \(F_x(u)=f(xu)\) on \(H\). The constants change only by the
fixed linear equivalence between \((a,b,z)\) and the standard Heisenberg basis. Thus
\[
\left(\int_0^\infty
\frac{\left(\int_H \abs{f(xu\exp(tz))-f(xu)}\dd u\right)^4}{t^3}\dd t\right)^{1/4}
\leq
C_{a,b}\int_H \bigl(\abs{a f(xu)}+\abs{b f(xu)}\bigr)\dd u .
\]
Integrating over \(G/H\) and using Minkowski's integral inequality yields
\[
\left(\int_0^\infty \frac{A_zf(t)^4}{t^3}\dd t\right)^{1/4}
\leq
C_{a,b}\int_G \bigl(\abs{a f(g)}+\abs{b f(g)}\bigr)\dd g .
\]
The final inequality follows because \(a,b\in V_1\) are fixed horizontal vectors, so their
directional derivatives are bounded by constants times \(\abs{\nabla_H f}\).
\end{proof}

We now prove Theorem \ref{thm:main}. Since the bracket map
\(\Lambda^2 V_1\to V_2\) spans \(V_2\), choose \(a_j,b_j\in V_1\), \(1\leq j\leq m\), such that
\[
v=\sum_{j=1}^m z_j,\qquad z_j=[a_j,b_j].
\]
Terms with \(z_j=0\) may be omitted. Since \(V_2\) is central,
\[
\exp(tv)=\prod_{j=1}^m \exp(tz_j).
\]
Let \(w_0=0\) and \(w_j=z_1+\cdots+z_j\). For each \(t>0\),
\begin{align*}
A_vf(t)
&=\int_G \abs{f(g\exp(tw_m))-f(g)}\dd g\\
&\leq \sum_{j=1}^m
\int_G
\abs{f(g\exp(tw_j))-f(g\exp(tw_{j-1}))}\dd g\\
&=\sum_{j=1}^m A_{z_j}f(t),
\end{align*}
where the last equality uses right-translation invariance of Haar measure by the central element
\(\exp(tw_{j-1})\). Applying Minkowski in \(L^4((0,\infty),t^{-3}\dd t)\) and then
Lemma \ref{lem:elementary},
\[
\left(\int_0^\infty \frac{A_vf(t)^4}{t^3}\dd t\right)^{1/4}
\leq
\sum_{j=1}^m
\left(\int_0^\infty \frac{A_{z_j}f(t)^4}{t^3}\dd t\right)^{1/4}
\leq C_{G,v}\int_G \abs{\nabla_H f}\dd g .
\]
This proves \eqref{eq:q4}.

For every \(q\geq 4\), the \(q=\infty\) deterministic path bound gives
\(\sup_{t>0} A_vf(t)/\sqrt t\lesssim_{G,v}\int_G\abs{\nabla_H f}\). Combining this with
\eqref{eq:q4},
\[
\int_0^\infty \left(\frac{A_vf(t)}{\sqrt t}\right)^q\frac{\dd t}{t}
\leq
\left(\sup_{t>0}\frac{A_vf(t)}{\sqrt t}\right)^{q-4}
\int_0^\infty \left(\frac{A_vf(t)}{\sqrt t}\right)^4\frac{\dd t}{t},
\]
so the endpoint inequality also holds for all \(q\geq 4\).

\section*{Verification notes and limitations}

The proof is exact for the continuous step-two Carnot subcase of Question 38 / Question 42.
It is not an optimal-exponent theorem: for \(\HH^{2k+1}\), \(k\geq 2\), Ryoo records the sharper
known value \(q=2\), while the argument here gives only the weaker admissible exponent \(q=4\).
No claim is made for higher-step Carnot groups or for arbitrary non-Carnot nilpotent Lie groups.

The main transfer point for review is Lemma \ref{lem:elementary}: one should check the use of
Weil disintegration over left cosets of the elementary Heisenberg subgroup and the harmless
linear rescaling of Naor--Young's standard \(\HH^3\) theorem.

The novelty check was bounded. I searched the run indexes for arXiv:2207.11305 and for
``step-two Carnot'', ``vertical perimeter'', ``L1'', ``Naor'', ``Young'', and Heisenberg variants.
External web search for close phrases found the Naor--Young Heisenberg endpoint results but no
explicit general step-two Carnot extension. The result should therefore be treated as a likely new
partial observation, pending expert literature review.

\begin{thebibliography}{9}

\bibitem{Ryoo2025}
S.-Y. Ryoo.
\newblock Quantitative nonembeddability of groups of polynomial growth into uniformly convex spaces.
\newblock arXiv:2207.11305, version 2, 2025.

\bibitem{NaorYoung2022}
A. Naor and R. Young.
\newblock Foliated corona decompositions.
\newblock arXiv:2004.12522; Acta Mathematica, 229, 55--200, 2022.

\end{thebibliography}

\end{document}
